% header.tex  (Quarto -> include-in-header)
% Ziel:
% - Seitenzahl im Kopf (mittig)
% - auch auf TOC/Verzeichnissen sichtbar
% - auch auf Kapitelanfangsseiten sichtbar (z.B. 1 Einleitung)
% - keine Linien/Striche

\usepackage[automark]{scrlayer-scrpage}
\usepackage{etoolbox}
\usepackage{caption}
\captionsetup[table]{justification=centering, singlelinecheck=false}

% Alle Kopf-/Fußdefinitionen löschen
\clearpairofpagestyles

% Seitenzahl oben mittig
\chead{\pagemark}
\cfoot{} % unten leer

% Linien AUS (hart)
\KOMAoptions{headsepline=false,footsepline=false}
\setheadsepline{0pt}
\setfootsepline{0pt}

% Kapitelanfangsseiten: gleichen Stil erzwingen
% (bei scrreprt ist das entscheidend, sonst bleibt Seite 1 der Einleitung "ohne Zahl")
\renewcommand*\chapterpagestyle{scrheadings}

% Standard-Pagestyle setzen
\pagestyle{scrheadings}

% TOC / Tabellen- / Abbildungsverzeichnis:
% Erste Seite bekommt oft \thispagestyle{empty/plain} -> wir überschreiben das.
\pretocmd{\tableofcontents}{\clearpage\thispagestyle{scrheadings}}{}{}
\pretocmd{\listoftables}{\clearpage\thispagestyle{scrheadings}}{}{}
\pretocmd{\listoffigures}{\clearpage\thispagestyle{scrheadings}}{}{}

% Zusätzlich: falls irgendwo "plain" benutzt wird, mache plain = scrheadings
\makeatletter
\let\ps@plain\ps@scrheadings
\makeatother